\section{Introduction}\label{sec:introduction}

Let \(S\) be finite and let \(M=(m_{st})_{s,t\in S}\) be symmetric,
with \(m_{ss}=1\) and
\[
 m_{st}\in\{2,3,\ldots\}\cup\{\infty\}\qquad(s\ne t).
\]
The Artin group \(A_M\) has generators \(\sigma_s\), \(s\in S\).
For each finite \(m_{st}\), its defining relation equates the two
alternating words of length \(m_{st}\) starting with \(\sigma_s\) and
\(\sigma_t\).  A label \(\infty\) imposes no relation.

Throughout, an action is called \emph{geometric} if it is an isometric
action on a nonempty proper CAT(0) metric space \(X\), is proper, and is
cocompact.  Here a metric space is proper when its closed bounded subsets
are compact; action properness means that
\[
 \{g:gK\cap K\ne\varnothing\}
\]
is finite for every compact \(K\subset X\); and cocompactness means that
\(GK=X\) for some compact \(K\).  Our CAT(0) convention is the geodesic
triangle-comparison inequality with Euclidean comparison triangles.
The explicit properness assumption on \(X\) also follows from the other
action hypotheses for a length space
\cite[Exercise~I.8.4(1)]{BridsonHaefliger1999}.

Charney asked which Artin groups are CAT(0)
\cite[Problem~4]{CharneyProblems}.  The affirmative expectation that every
finite-rank Artin group admits a geometric action is recorded by Haettel
as Conjecture~1.1 of \cite{Haettel2022}.  A geometric action
would make the group's large-scale geometry accessible through geodesic
convexity and Euclidean comparison.  The challenge is compatibility:
constructing actions for the rank-two braid subgroups does not by itself
supply an action of the whole group.

There are substantial positive results.  Right-angled Artin groups,
whose finite off-diagonal labels are all two, have finite nonpositively
curved cubical classifying spaces and hence geometric CAT(0) actions
\cite[preprint Theorem~3.1.1]{CharneyDavis1995}.
Brady and McCammond constructed
nonpositively curved two-dimensional complexes for the three-generator
Artin groups with all three finite labels at least three
\cite[preprint Theorem~4]{BradyMcCammond2000}.  Haettel proved the conjecture for
XXL type, where every finite off-diagonal label is at least five
\cite[Theorem~1.4]{Haettel2022}.
For classical braid groups, Brady and McCammond established the
five-strand case \cite[Theorem~B]{BradyMcCammond2010},
Haettel, Kielak, and Schwer the six-strand case
\cite[preprint Corollary~4.18]{HaettelKielakSchwer2016}, and Jeong the
seven-strand case \cite[preprint Corollary~14]{Jeong2023}.

In the opposite direction, Brady and Crisp found Artin groups with
two-dimensional classifying spaces for which the least dimension of
a geometric CAT(0) action is three
\cite[preprint Theorems~B and~3.1]{BradyCrisp2002}.
Their examples show that an
obstruction in one dimension can disappear in a higher-dimensional
CAT(0) space.  To refute the conjecture, an obstruction must instead
survive every dimension.

\begin{theorem}\label{thm:main}
There is an Artin matrix \(M\) on 116 generators, with every
off-diagonal label in \(\{2,3,\infty\}\), such that \(A_M\) admits no
geometric action.  One explicit matrix is given in
Definition~\ref{def:matrix}.
\end{theorem}

Thus Theorem~\ref{thm:main} gives a negative resolution of the
CAT(0) conjecture for Artin groups.
In particular, the conclusion excludes every nonempty proper CAT(0)
space, regardless of dimension or the choice of a complex or metric.

\subsection*{The construction and the obstruction}

The small group that detects the obstruction is
\[
 T=\langle a,b,c\mid aba=bab,\ bcb=cbc,\ cac=aca\rangle.
\]
To each of its three cyclic braid pairs we attach a 40-strand braid
group, sharing the first two standard generators of that block with the
pair.  Each attachment adds 37 generators.  Two further generators
\(d,e\), each commuting with \(a,b,c\) and having no other finite
incident labels, make \(T\) their simultaneous centralizer;
Section~\ref{sec:assembly} proves this equality.  This gives \(3+3\cdot37+2=116\) generators.

For an isometry coming from a geometric action, write \(\ell(g)\) for its
stable translation length, the limit of \(d(x,g^n x)/n\).
Section~\ref{sec:lengths} turns these lengths on commuting elements into
Euclidean pairings and proves that a pairing with a fixed element is
additive on its whole centralizer.  The calculation in
Section~\ref{sec:braids} uses these standard geometric properties to
constrain the braid relations, without requiring the braid blocks
themselves to act cocompactly.  A long block forces
\[
 \frac{\ell((uv)^3)}{3\ell(u^2)}
 =\frac{\ell((uv)^3)}{3\ell(v^2)}>\frac12.
\]
If these inequalities held on all three pairs in a geometric action of
\(T\), the commuting pairs
\[
 (a^2,(ab)^3),\ (b^2,(ab)^3),\
 (b^2,(bc)^3),\ (c^2,(bc)^3),\
 (c^2,(ca)^3),\ (a^2,(ca)^3)
\]
would give six asymptotically Euclidean sectors, each of angle less than
\(\pi/3\).  Their cyclic union would have total angle less than \(2\pi\).

A related angle-sum obstruction appears in Crisp's study of
two-dimensional Bestvina--Brady groups
\cite[Lemma~1.3 and the following remark]{Crisp2002}.
His argument assumes covering dimension two and a proper action in which
each isometry attains its minimum displacement.

The main geometric step is to turn this short angular circuit into a
contradiction that remains valid in arbitrary dimension.  Coordinates
coming from an affine mirror action separate disjoint regions of the
sector disk at large scale.  This action is the equilateral-triangle
specialization of Digne's reflection-coordinate homomorphism
\cite[Proposition~2.1]{Digne2016}.  CAT(0) geodesics provide a second
disk whose image stays far from the cone point.  We transfer the resulting sphere
to a triangular presentation complex, using a bound valid at every
point of every parameter triangle.  A local winding count then detects
a target triangle with nonzero signed multiplicity, whereas a direct
folding argument makes every such multiplicity zero.

The triangular presentation and its girth-six link are the
\(m=n=p=3\) specialization of Brady and McCammond's
construction~\cite[preprint Lemmas~2--3]{BradyMcCammond2000}.
We prove the facts needed here,
including the signed-count and controlled-approximation statements.
The latter comparison between local degree and global signed counts
is formulated entirely in terms of distances and finite diagrams.
It does not require a triangulation or a dimension bound on the unknown
CAT(0) space.

Finally, the simultaneous-centralizer construction makes the action on
\(T\) geometric: a suitable intersection of two displacement sublevel
sets is closed and convex, and \(T\) acts cocompactly on it.  Its stable
lengths are the ambient stable lengths.  The braid inequalities and the
six-sector obstruction therefore apply to the same action.

\subsection*{Organization and notation}

Sections~\ref{sec:lengths} and~\ref{sec:braids} establish the length
calculation.  Sections~\ref{sec:sectors}--\ref{sec:localization} prove the
obstruction for \(T\).  Section~\ref{sec:assembly} gives the matrix,
the algebraic embeddings, and the centralizer argument completing
Theorem~\ref{thm:main}.  The substantial geometric, algebraic, and
diagrammatic ingredients are proved below.  We use elementary finite planar topology for regular
neighbourhoods of graphs and winding numbers; no Artin-group asphericity
theorem is an input.

We use \(B_l\) for the braid presentation on
\(s_1,\ldots,s_{l-1}\): adjacent generators satisfy the length-three
braid relation, and generators at distance at least two commute.
For commuting elements \(g,h\), the notation \(\langle g,h\rangle\)
will denote the polarization of their squared stable lengths, defined
in Lemma~\ref{lem:commuting-form}.
